% =====================================================================
% APPENDIX B (PRDL, Maria Martinez): learning strategies, main takeaways
% Asked for in the weekly summary of 9/10 (quotes.md F5); material from
% docs/literature/learning_strategies.md (T18), checked quotes only.
% =====================================================================
\appendixbanner{APPENDIX B: LEARNING STRATEGIES, MAIN TAKEAWAYS}
\addcontentsline{toc}{section}{Appendix B: Learning strategies, main takeaways}\label{app:strategies}

The weekly summary of 9 October asks to ``include an annex with your main takeaways'' on learning
strategies in the final PRDL report \cite{course_summary_0910}. The class presentations of the first
assignment are not available to us, so this annex is written from checked quotes and from what area~2 of
our project did (Section~\ref{sec:strategies}).

\subsection*{1. Transfer learning}

Transfer learning reuses what was learnt on one task or data distribution for another; a survey defines
it as the area that ``groups approaches dealing with such an unpredictable distribution shift''
\cite[p.~2]{monka2022survey}. In the course view it usually relies on manually labelled data in the source
task \cite{course_summary_0910}. A pre-trained network can be reused as a frozen feature extractor with a
new linear classifier, or fine-tuned end to end \cite[p.~11]{radford2021learning}; fine-tuning uses a smaller
learning rate so as not to overwrite what was learnt, ``to prevent catastrophic forgetting''
\cite[p.~6]{soltanieh2023indistribution}. On ECG, the benefit of pre-training grows as the target training set
shrinks \cite[p.~1524]{strodthoff2021deep}.

\emph{Takeaway from our project.} The frozen ECG-GPT with a linear probe was weaker than training from
scratch with all the labels; fine-tuning was better than both, with all the labels and with 10\,\% of
them. ``Transfer'' is not one thing: what is transferred (frozen features or a starting point) decides the
result.

\subsection*{2. Self-supervised learning}

Self-supervised learning learns features without manual labels, to overcome their scarcity
\cite[p.~1]{soltanieh2022analysis}. The model solves a pretext task built from the data itself and is then reused;
the course notes that transfer and self-supervised learning both reuse knowledge and differ in the labels
used for pre-training \cite{course_summary_0910}. On 12-lead ECG, a self-supervised model trained on half
of the labelled data approximately reached the supervised performance at twice that amount
\cite[pp.~7--8]{mehari2022self}. A generative pretext task is the GPT recipe: a decoder-only Transformer trained to
predict the next token \cite[p.~3]{radford2018improving}, an idea that has also been applied to ECG tokens
\cite[p.~1]{yang2025heartllm}.

\emph{Takeaway from our project.} Next-sample prediction on lead II, without labels, gave a starting point
that made the classifier far more label-efficient. But the generator chosen for being the best generator
gave the worst frozen features: a good pretext model for one task is not automatically the best
representation for another.

\subsection*{3. Zero-shot learning}

In zero-shot learning the classes of the target task are not seen with labels during training; in a
visual-transfer survey it is defined as a task ``with labeled source domain data and unlabeled target
domain data'' \cite[p.~4]{monka2022survey}. For NLP, the course takes its definitions of zero-shot and few-shot from the paper that
presented GPT-3 \cite{course_summary_0910}. With other modalities,
zero-shot needs the inputs and the class descriptions to share an embedding space (``multimodal
correspondence'') \cite{course_summary_0910}. CLIP learns such a space from image--text pairs collected on
the Internet and classifies by the similarity between an image and a sentence for each class
\cite[pp.~3, 6]{radford2021learning}.

\emph{Takeaway from our project.} CLIP worked zero-shot on everyday images (CIFAR-100) and recognised that
our images were electrocardiograms, but it could not diagnose them. Zero-shot works only as far as the
shared space covers the task. Why it does not cover ECG diagnosis is a hypothesis: the source does not say
whether ECG traces with diagnoses are among CLIP's training pairs.

\subsection*{4. Connections}

The three strategies are not exclusive. Our ECG-GPT is self-supervised pre-training followed by transfer;
CLIP is pre-training on paired data followed by zero-shot use. Which one fits a problem depends on what is
available: labels for the target task (transfer and fine-tuning), unlabelled signals (self-supervised
pre-training), or a space shared with text (zero-shot). The fewer the target labels, the more
pre-training is worth, which is why we ran the transfer experiment with 100\,\% and 10\,\% of the labels.
